\chapter{Introduction}
\label{ch:intro}

\section{Objectif du rapport}

Ce rapport documente l'utilisation de l'intelligence artificielle dans le projet
\textbf{Promptly}. Il presente :
\begin{itemize}
  \item les \textbf{outils d'IA} utilises pour concevoir et coder l'application ;
  \item les \textbf{fonctionnalites d'IA} integrees au produit livre.
\end{itemize}

\section{Deux faces de l'IA}

\begin{enumerate}
  \item \textbf{IA cote developpement.} ChatGPT a servi a preparer le premier
        grand prompt. L'editeur \textbf{Cursor} (mode Agent) a ensuite heberge
        le dialogue avec des \textbf{modeles} tels que \textbf{Claude} et
        \textbf{Grok}, qui ont genere, corrige et explique le code. Chaque etape
        importante est documentee : prompt, resultat, eventuel probleme, prompt
        de correction, nouveau resultat --- avec captures d'ecran.
  \item \textbf{IA cote application.} Le produit appelle des modeles et des APIs
        (Ollama, Hugging Face, Sightengine, TextGears, Gemini, services ML locaux)
        pour resume, chat, moderation, emotion, matching CV, etc.
\end{enumerate}

\section{Methode de documentation}

Pour chaque echange significatif avec une IA de developpement :
\begin{enumerate}
  \item capture du \textbf{prompt} ;
  \item capture du \textbf{resultat} ;
  \item si erreur ou ecart : capture du \textbf{probleme}, du \textbf{prompt de
        correction}, puis du \textbf{resultat corrige}.
\end{enumerate}

Les images sont placees dans \texttt{figures/chatgpt/}, \texttt{figures/cursor/}
et \texttt{figures/app/}.

\section{Organisation du document}

Le chapitre~\ref{ch:role} rappelle brievement le role de Promptly.
Le chapitre~\ref{ch:dev} detaille la chaine ChatGPT~$\rightarrow$~Cursor
(Claude / Grok).
Le chapitre~\ref{ch:app} presente l'IA integree au produit.
Le chapitre~\ref{ch:conclusion} conclut.

% --- Contenu a completer apres validation ---
